\section{Five-vertex branch-DAG encoding and checker contract}
\label{app:n5-dag}

The five-vertex lower certificates prove infeasibility of a hitting-set budget
on an explicit core of pair-separator masks.  This appendix specifies why a
locally accepted DAG certifies the global lower bound.

\subsection{Root data}

For each action the proof object contains the ordered coordinate IDs, the
budget, and a list of explicit unordered pattern pairs.  The checker rebuilds
each canonical pattern from its ID and recomputes the integer coordinate rows
before deriving the separator mask.  The root statistics are:
\[
\begin{array}{c|r|r|l}
G & \text{budget} & \text{core rows} & \text{root-state SHA-256 prefix}\\ \hline
S_5 & 13 & 249 & \texttt{80a7d31c66c96555}\\
A_5 & 16 & 285 & \texttt{0fb0e501622ad7bc}.
\end{array}
\]
The pair core is a subset of the complete separator system.  Therefore, if no
budget-$b$ family hits the core, no budget-$b$ family separates all pattern
orbits.

\subsection{Reduced states and local rules}

A state consists of a nonnegative remaining budget and a set of uncovered
separator masks.  Selecting coordinate $i$ removes every mask containing
$i$ and decreases the budget by one.  Singleton masks are propagated until
none remains.  The state is rejected if the budget is negative or a remaining
mask is empty.

A \emph{branch node} names one uncovered mask $S$ and contains one child for
every $i\in S$, after selecting $i$ and applying deterministic singleton
propagation.  Since a hitting set must meet $S$, these children exhaust all
solutions below that node.

A \emph{packing leaf} supplies $k$ pairwise disjoint uncovered masks when
$k$ exceeds the remaining budget.  Each selected coordinate hits at most one
packing member, so closure is valid.  A \emph{coverage leaf} supplies the
number $C$ of uncovered masks and an exact maximum $M$ of masks hit by any
remaining coordinate.  If $\lceil C/M\rceil$ exceeds the remaining budget,
closure is again valid.

Every state stores a canonical state hash, but the checker recomputes the
reduced masks, budgets, packing disjointness, coverage maximum, and child
states.  The hash binds the representation after those semantic checks.

\begin{lemma}[Soundness of a complete branch DAG]
If every reachable node satisfies its local transition rule and every leaf
satisfies one of the two lower-bound rules above, then the root state has no
hitting set within its budget.
\end{lemma}

\begin{proof}
Proceed by reverse topological order.  A checked leaf has no feasible hitting
set by its packing or coverage inequality.  At a branch node, any hitting set
must select some coordinate in the named mask and would therefore induce a
feasible hitting set in at least one child.  If all children are infeasible,
so is the parent.  The conclusion reaches the root.  Shared child states do
not affect the induction because the graph is acyclic and every reachable
node is checked once.
\end{proof}

\subsection{Complete proof census}

For $S_5$ there are $11{,}893$ branch nodes and $51{,}696$ leaves.  Of the
leaves, $49{,}123$ close by disjoint packing and $2573$ by uniform coverage.
The maximum depth is eleven.  For $A_5$ there are $44{,}316$ branch nodes and
$169{,}332$ leaves: $164{,}003$ packing leaves and $5329$ coverage leaves,
with maximum depth fourteen.  All $63{,}589$ and $213{,}648$ nodes are
reachable from their roots.

The deterministic compressed transports and canonical uncompressed content
are separately hashed:
\begin{center}
\small
\begin{tabular}{ccl}
\toprule
Action & representation & SHA-256\\
\midrule
$S_5$ & gzip & \texttt{6babe9631d4dc613\ldots f8c60d}\\
$S_5$ & canonical JSON & \texttt{233be7b8bb5aa9ef\ldots200a2d}\\
$A_5$ & gzip & \texttt{11d3efdc8b683965\ldots d2cde}\\
$A_5$ & canonical JSON & \texttt{866974b4163f90e3\ldots eab82}\\
\bottomrule
\end{tabular}
\end{center}
The full hashes are recorded in the atlas manifest.  The verifier checks both
transport forms, but the lower-bound conclusion follows from the local rules,
not from matching the displayed prefixes.
